% GENERATED FILE. Edit canonical lessons, then run npm run generate:books.
% canonical-source-hash: fnv1a64:2a8e2154692530c8
% canonical-lessons: GE-C281-einbrechen, GE-C281-einbruch, GE-C281-einbrecher, GE-C281-klauen, GE-C281-beute

\chapter{Der Einbruch}
\label{ch:ge-der-einbruch}
\hlchaptermodality{281}

\begin{chapteropening}
\textbf{By the end of this chapter:} \emph{I can report a break-in: to break in, a burglary, a burglar, to pinch and the loot.}

Complete the last lesson of chapter 281: I can report a break-in: to break in, a burglary, a burglar, to pinch and the loot.
\end{chapteropening}

\section[einbrechen]{einbrechen — to break in, to burgle}

\label{lesson:GE-C281-einbrechen}



\begin{quote}
Before the new one: say the German for the victim, then the German for the crime scene.
\end{quote}

\subsection*{You'll want to know: einbrechen}
\textbf{einbrechen} — \textquotedblleft{}to break in, to burgle\textquotedblright{}.

\begin{quote}
In der Nacht brachen Diebe in das Geschäft ein.
\end{quote}

\begin{quote}
\emph{At night thieves broke into the shop.}
\end{quote}

Separable, Perfekt with \emph{sein}: \emph{Jemand ist eingebrochen}. Also of ice: \emph{Das Eis brach ein}.

\begin{cousinweb}
\emph{einbrechen} joins \emph{ein} \textquotedblleft{}in\textquotedblright{} and \emph{brechen} \textquotedblleft{}to break\textquotedblright{}, from Old High German \emph{brehhan}, Proto-Germanic \emph{brekaną} (EtymDB).
\end{cousinweb}

\begin{grammarlens}[title={What you've built}]
One more word for telling, explaining and describing at B1.
\end{grammarlens}

\subsection*{Your turn}
\begin{itemize}
\raggedright
  \item \emph{Say it:} \emph{einbrechen}
  \item \emph{Say it:} \emph{einbrechen}, once more
  \item \emph{Say it:} a sentence of your own with \emph{einbrechen}
  \item \emph{Recall:} say the German for the victim, then the German for the crime scene, then say \emph{einbrechen} again
\end{itemize}

\begin{tcolorbox}[breakable,colback=teal!4,colframe=teal!35!black,title={Before you move on}]
What does einbrechen mean? (\textquotedblleft{}To break in, to burgle\textquotedblright{}.) Say it once more.
\end{tcolorbox}
\section[der Einbruch]{der Einbruch — the burglary, the break-in}

\label{lesson:GE-C281-einbruch}



\begin{quote}
Before the new one: say the German for the crime scene, then the German for to break in.
\end{quote}

\subsection*{You'll want to know: der Einbruch}
\textbf{der Einbruch} — \textquotedblleft{}the burglary, the break-in\textquotedblright{}.

\begin{quote}
Nach dem Einbruch fehlte der Fernseher.
\end{quote}

\begin{quote}
\emph{After the burglary the television was missing.}
\end{quote}

Plural: \emph{die Einbrüche}. \emph{Bei Einbruch der Nacht} — at nightfall.

\begin{cousinweb}
\emph{Einbruch} is built on \emph{einbrechen} \textquotedblleft{}to break in\textquotedblright{}, from \emph{brechen} \textquotedblleft{}to break\textquotedblright{}, Proto-Germanic \emph{brekaną} (EtymDB).
\end{cousinweb}

\begin{grammarlens}[title={What you've built}]
One more word for telling, explaining and describing at B1.
\end{grammarlens}

\subsection*{Your turn}
\begin{itemize}
\raggedright
  \item \emph{Say it:} \emph{der Einbruch}
  \item \emph{Say it:} \emph{der Einbruch}, once more
  \item \emph{Say it:} a sentence of your own with \emph{der Einbruch}
  \item \emph{Recall:} say the German for the crime scene, then the German for to break in, then say \emph{der Einbruch} again
\end{itemize}

\begin{tcolorbox}[breakable,colback=teal!4,colframe=teal!35!black,title={Before you move on}]
What does der Einbruch mean? (\textquotedblleft{}The burglary, the break-in\textquotedblright{}.) Say it once more.
\end{tcolorbox}
\section[der Einbrecher]{der Einbrecher — the burglar}

\label{lesson:GE-C281-einbrecher}



\begin{quote}
Before the new one: say the German for to break in, then the German for the burglary.
\end{quote}

\subsection*{You'll want to know: der Einbrecher}
\textbf{der Einbrecher} — \textquotedblleft{}the burglar\textquotedblright{}.

\begin{quote}
Der Einbrecher kam durch das Fenster.
\end{quote}

\begin{quote}
\emph{The burglar came through the window.}
\end{quote}

Plural: \emph{die Einbrecher}. Feminine \emph{die Einbrecherin}.

\begin{cousinweb}
\emph{Einbrecher} is built on \emph{einbrechen} \textquotedblleft{}to break in\textquotedblright{}, from \emph{brechen} \textquotedblleft{}to break\textquotedblright{} (EtymDB), with the ending \emph{-er} for a person.
\end{cousinweb}

\begin{grammarlens}[title={What you've built}]
One more word for telling, explaining and describing at B1.
\end{grammarlens}

\subsection*{Your turn}
\begin{itemize}
\raggedright
  \item \emph{Say it:} \emph{der Einbrecher}
  \item \emph{Say it:} \emph{der Einbrecher}, once more
  \item \emph{Say it:} a sentence of your own with \emph{der Einbrecher}
  \item \emph{Recall:} say the German for to break in, then the German for the burglary, then say \emph{der Einbrecher} again
\end{itemize}

\begin{tcolorbox}[breakable,colback=teal!4,colframe=teal!35!black,title={Before you move on}]
What does der Einbrecher mean? (\textquotedblleft{}The burglar\textquotedblright{}.) Say it once more.
\end{tcolorbox}
\section[klauen]{klauen — to pinch, to nick (colloquial)}

\label{lesson:GE-C281-klauen}



\begin{quote}
Before the new one: say the German for the burglary, then the German for the burglar.
\end{quote}

\subsection*{You'll want to know: klauen}
\textbf{klauen} — \textquotedblleft{}to pinch, to nick (colloquial)\textquotedblright{}.

\begin{quote}
Jemand hat mein Fahrrad geklaut.
\end{quote}

\begin{quote}
\emph{Someone has nicked my bike.}
\end{quote}

Colloquial for \emph{stehlen}. Regular: \emph{klaute}, \emph{hat geklaut}.

\begin{cousinweb}
\emph{klauen} is usually linked to \emph{Klaue} \textquotedblleft{}claw\textquotedblright{}, from an older verb meaning \textquotedblleft{}to grab with the claws\textquotedblright{}; EtymDB lists no parent, so the link is not settled.
\end{cousinweb}

\begin{grammarlens}[title={What you've built}]
One more word for telling, explaining and describing at B1.
\end{grammarlens}

\subsection*{Your turn}
\begin{itemize}
\raggedright
  \item \emph{Say it:} \emph{klauen}
  \item \emph{Say it:} \emph{klauen}, once more
  \item \emph{Say it:} a sentence of your own with \emph{klauen}
  \item \emph{Recall:} say the German for the burglary, then the German for the burglar, then say \emph{klauen} again
\end{itemize}

\begin{tcolorbox}[breakable,colback=teal!4,colframe=teal!35!black,title={Before you move on}]
What does klauen mean? (\textquotedblleft{}To pinch, to nick (colloquial)\textquotedblright{}.) Say it once more.
\end{tcolorbox}
\section[die Beute]{die Beute — the loot; the prey}

\label{lesson:GE-C281-beute}



\begin{quote}
Before the new one: say the German for the burglar, then the German for to pinch.
\end{quote}

\subsection*{You'll want to know: die Beute}
\textbf{die Beute} — \textquotedblleft{}the loot; the prey\textquotedblright{}.

\begin{quote}
Die Diebe teilten die Beute.
\end{quote}

\begin{quote}
\emph{The thieves shared out the loot.}
\end{quote}

No plural in this sense. For animals, the prey.

\begin{cousinweb}
\emph{Beute} comes from Middle High German \emph{biute}, Old High German \emph{biuta} (EtymDB), usually traced to a Low German word for \textquotedblleft{}exchange, sharing out\textquotedblright{}.
\end{cousinweb}

\begin{grammarlens}[title={What you've built}]
One more word for telling, explaining and describing at B1.
\end{grammarlens}

\subsection*{Your turn}
\begin{itemize}
\raggedright
  \item \emph{Say it:} \emph{die Beute}
  \item \emph{Say it:} \emph{die Beute}, once more
  \item \emph{Say it:} a sentence of your own with \emph{die Beute}
  \item \emph{Recall:} say the German for the burglar, then the German for to pinch, then say \emph{die Beute} again
\end{itemize}

\begin{tcolorbox}[breakable,colback=teal!4,colframe=teal!35!black,title={Before you move on}]
What does die Beute mean? (\textquotedblleft{}The loot; the prey\textquotedblright{}.) Say it once more.
\end{tcolorbox}
